% ===========================================================================
\section{Introduction}
\label{sec:introduction}
% ===========================================================================

Household fuel is the sector where greenhouse-gas accounting and air-quality accounting most
clearly disagree about what matters.

Under the IPCC conventions, burning dung and firewood is a modest source: the \COtwo{} is
biogenic and reported outside the national total, leaving only the methane and nitrous oxide
of incomplete combustion. Under a primary-emission lens the same combustion is among the dominant sources of
primary particulate in the country, emitted at breathing height inside and around
dwellings by the population least able to avoid it---though translating emitted mass
into exposure requires evidence this inventory does not carry
(Section~\ref{sec:limitations}).

Both accounts are correct. They are answering different questions, and a sector that is
\SI{5.4}{} or \SI{6.5}{\percent} of Pakistan's emissions (under the survey-based reference and
the all-diverted dung configuration of Section~\ref{sec:results}) and, under either, one of
the largest sources of its primary particulate is where the difference between those
questions has the largest practical consequence.

\subsection{What this paper does}

We report Pakistan's residential fuel combustion for FY2023--24 from national fuel supply
statistics and household-survey quantities, with an air-pollutant layer computed from the same fuel quantities, and treat the
comparison between the two as the result rather than as an aside. Throughout,
``layers'' distinguishes the greenhouse-gas from the air-pollutant calculation,
``configurations'' the two dung activity bases (S, the reference, and L), and ``scenarios'' the
structural sensitivity cases; the word inventory refers to the whole system.

The comparison is defensible here in a way it usually is not. Cross-source particulate
comparisons in developing countries are typically assembled from separate studies with
separate activity bases, so the ratio between two sources is partly an artefact of which
datasets each author had. Every air layer in the inventory this sector belongs to is computed
from that sector's own greenhouse-gas activity data---the same fuel, the same tonnage, the
same population panel---so the ratio between household fuel and road transport
reported in Section~\ref{sec:air} (\numrange{49}{109} across the two dung
configurations) is free of the mismatched-activity artefact, though it remains
conditional on the stated factor, boundary and configuration choices.

\subsection{A transfer that had nowhere to go}

The sector also makes explicit a gap that had been open in this inventory and is common in others.
Pakistani households burn cattle and buffalo dung, and the agriculture sector accounts for
that on the supply side: it removes \SI{28}{\percent} of cattle and \SI{30}{\percent} of
buffalo excreted nitrogen from the manure-management balance and reports it with the notation
key IE---included elsewhere.

Until a residential sector existed, nowhere was elsewhere: \SI{1.259}{\mega\tonne} of
excreted nitrogen a year left one balance correctly and entered no other. The ledger now
conserves that diverted quantity---what agriculture removes is accounted for exactly once,
as household combustion or as a typed unresolved remainder---while the \emph{destination}
of the diverted mass is treated as an
open question: the reference books to household combustion the consumption the household
survey observes, leaving \SI{47.8}{\mega\tonne} with no observed destination, outside the
total; configuration L, the reference of earlier releases, books it all to household
combustion
(Section~\ref{sec:methods}).

\subsection{Contributions}

\begin{enumerate}[leftmargin=*,itemsep=2pt,topsep=3pt]
\item \textbf{A fuel-anchored household energy inventory for Pakistan} across nine streams,
  with no household-count model and therefore nothing that can drift from the fuel actually
  supplied (Section~\ref{sec:methods}).
\item \textbf{An executable cross-sector dung ledger} that conserves the quantity
  diverted from manure management while distinguishing the observed household
  destination from an unresolved remainder (\ref{app:dung}).
\item \textbf{Greenhouse gases and air pollutants on one activity basis}, which is what makes
  the source ranking in Section~\ref{sec:air} a physical statement rather than a
  data-availability artefact.
\item \textbf{An explicit treatment of the biomass memo convention}, under which
  \SI{65}{\percent} of household energy produces \SI{24}{\percent} of reported emissions, and
  of the \SI{65}{\mega\tonne}-scale biogenic \COtwo{} that sits outside the total.
\item \textbf{A survey-anchored firewood quantity}---the HIES 2024--25 design-weighted
  household measurement, corroborated by the independent 2018--19 round and by FAOSTAT's
  moisture-matched equivalent, replacing the flagged FAO estimate the sector previously
  rode---with the sensitivity still published rather than a single number
  (Section~\ref{sec:results}).
\end{enumerate}

Method shared with the other nine sectors is documented separately
\citep{ahmad2026national}.
